\documentclass[11pt,a4paper]{article}
\usepackage{amsmath}
\usepackage{fontspec}
\usepackage{babel}
\usepackage[margin=2.3cm]{geometry}
\usepackage{enumitem}
\usepackage{xcolor}
\usepackage{hyperref}
\usepackage{mdframed}
\usepackage{booktabs}
\usepackage{array}
\usepackage{longtable}

\definecolor{primary}{HTML}{cc785c}
\definecolor{coral}{HTML}{cc785c}
\definecolor{ink}{HTML}{141413}
\definecolor{body}{HTML}{3d3d3a}
\definecolor{muted}{HTML}{6c6a64}
\definecolor{canvas}{HTML}{faf9f5}
\definecolor{surface-soft}{HTML}{f5f0e8}
\definecolor{hairline}{HTML}{e6dfd8}
\definecolor{accent-teal}{HTML}{5db8a6}
\definecolor{accent-amber}{HTML}{e8a55a}
\definecolor{lightbg}{HTML}{f5f0e8}

\hypersetup{colorlinks=true, linkcolor=coral, urlcolor=coral}

\babelprovide[import, onchar=ids fonts]{thai}
\babelprovide[import, onchar=ids fonts]{english}

\babelfont{rm}{Noto Sans}
\babelfont[thai]{rm}{Noto Sans Thai}
\babelfont{tt}{Noto Sans Mono}
\babelfont[thai]{tt}{Noto Sans Thai}

\directlua{
  local glyph_id = node.id("glyph")
  local penalty_id = node.id("penalty")
  local function is_thai(c) return c >= 0x0E01 and c <= 0x0E5B end
  local function is_thai_combining(c)
    return (c >= 0x0E31 and c <= 0x0E3A) or (c >= 0x0E47 and c <= 0x0E4E)
  end
  luatexbase.add_to_callback("pre_linebreak_filter", function(head)
    for n in node.traverse_id(glyph_id, head) do
      local prev = n.prev
      if prev and prev.id == glyph_id and is_thai(prev.char)
         and is_thai(n.char) and not is_thai_combining(n.char) then
        local p = node.new(penalty_id)
        p.penalty = 0
        node.insert_before(head, n, p)
      end
    end
    return head
  end, "thai_break")
}

\emergencystretch=1em
\setlist[itemize]{leftmargin=1.8em}
\setlist[enumerate]{leftmargin=2em}
\newcolumntype{L}[1]{>{\raggedright\arraybackslash}p{#1}}

\newmdenv[
  backgroundcolor=lightbg,
  linecolor=coral,
  linewidth=2pt,
  topline=false, rightline=false, bottomline=false,
  innerleftmargin=12pt, innerrightmargin=8pt,
  innertopmargin=8pt, innerbottommargin=8pt,
]{highlight}

\title{\textbf{รายละเอียด Paper ปี 2025 สำหรับ Research Gap}\\
\large Scientific Figure Captioning และ Context-Aware Caption Generation}
\author{Thesis Notes}
\date{24 พฤษภาคม 2026}

\begin{document}
\maketitle

\begin{highlight}
\textbf{โจทย์วิทยานิพนธ์:} การสร้างคำอธิบายรูปภาพและกราฟในเอกสารวิชาการโดยใช้บริบทจากเนื้อหาเอกสาร \\
\textbf{แกน keyword:} Scientific Figure Captioning, Context-Aware Caption Generation, Multimodal Profile, Figure-Mentioning Paragraph, OCR, Faithfulness Evaluation
\end{highlight}

\section*{ภาพรวม Paper ที่เลือก}

\begin{longtable}{@{}L{4.2cm}L{2cm}L{4.2cm}L{4.8cm}@{}}
\toprule
\textbf{Paper} & \textbf{ปี} & \textbf{แกนหลัก} & \textbf{เหตุผลที่เกี่ยวข้องกับ Thesis} \\
\midrule
LaMP-Cap: Personalized Figure Caption Generation With Multimodal Figure Profiles & 2025 & Dataset + multimodal profile & พิสูจน์ว่า context จาก paper เดียวกันช่วย caption generation และเปิดช่องให้ thesis ขยายเป็น full-document retrieval \\
Leveraging Author-Specific Context for Scientific Figure Caption Generation & 2025 & Two-stage context-aware prompt pipeline & แสดงปัญหา zero-shot hallucination และวิธี filter/context refinement ที่ใช้ได้จริงใน SciCap Challenge \\
Personalized Scientific Figure Caption Generation & 2025 & Fine-tuned VLM + quality-aware generation & ชี้ว่า paragraph, mention, OCR และ profile data ช่วย performance แต่เกิด trade-off ระหว่าง style กับ quality \\
\bottomrule
\end{longtable}

\newpage

\section{LaMP-Cap: Personalized Figure Caption Generation With Multimodal Figure Profiles}

\noindent
\begin{tabular}{@{}ll@{}}
\textbf{arXiv} & \href{https://arxiv.org/abs/2506.06561}{2506.06561} \\
\textbf{ผู้เขียน} & Ho Yin ``Sam'' Ng, Ting-Yao Hsu, Aashish Anantha Ramakrishnan และคณะ \\
\textbf{Keyword} & Scientific Figure Captioning, Personalized Caption Generation, Multimodal Profile \\
\textbf{ไฟล์ใน Repo} & papers/8-LaMPCap\_2025.* \\
\end{tabular}

\subsection*{ปัญหาที่ Paper แก้}
งาน image captioning หรือ figure captioning เดิมมักสร้าง caption จาก target image และข้อความประกอบเพียงเล็กน้อย แต่ scientific figure caption ใน paper จริงไม่ได้ต้องการแค่คำบรรยายภาพทั่วไป ผู้เขียนมักต้องการ caption ที่สอดคล้องกับ domain vocabulary, writing style และ narrative ของ paper เดียวกัน งาน personalization ก่อนหน้าอย่าง LaMP ส่วนใหญ่เน้น text-only setting จึงยังไม่ตอบโจทย์ multimodal scientific captioning ที่ทั้ง input และ profile มีทั้งภาพและข้อความ

\subsection*{Contribution สำคัญ}
\begin{enumerate}
\item เสนอ \textbf{LaMP-Cap dataset} สำหรับ personalized figure caption generation โดยสร้างจาก SciCap Challenge Dataset
\item นิยาม task เป็นการสร้าง caption ให้ target figure โดยใช้ multimodal source และ multimodal profile จาก paper เดียวกัน
\item ให้ profile figure สูงสุด 3 รูปต่อ target figure โดยแต่ละ profile มี image, caption และ figure-mentioning paragraph
\item ทดลองกับ LLM/VLM หลายตัวและพบว่า profile information ช่วยให้ generated captions ใกล้ ground-truth author captions มากขึ้น
\item ablation study แสดงลำดับความสำคัญของ profile component: caption สำคัญที่สุด ตามด้วย image และ paragraph
\end{enumerate}

\subsection*{Dataset และ Method}
LaMP-Cap curate จาก SciCap Challenge Dataset ซึ่งมี scientific figures จาก arXiv papers จำนวนมาก วิธีสร้างคือเลือกหนึ่ง figure จาก paper เป็น target figure แล้วใช้ figure อื่นจาก paper เดียวกันเป็น profile เพื่อแทน context และ author style เอกสารที่มีเพียงรูปเดียวถูกตัดออกเพราะไม่มี profile ให้ใช้

\begin{highlight}
\textbf{ตัวเลขสำคัญของ dataset}
\begin{itemize}[nosep]
\item 110,828 target figures
\item 307,903 figures รวม target และ profile
\item 197,075 profile figures
\item split: 80\% train, 10\% validation, 10\% test
\item figure types หลัก: Graph Plot, Node Diagram, Equation, Bar Chart, Scatterplot
\end{itemize}
\end{highlight}

Input ของ task แบ่งเป็นสองส่วน ส่วนแรกคือ target source ได้แก่ target figure image และ figure-mentioning paragraph ส่วนที่สองคือ profile ได้แก่ image, caption และ related paragraph ของ figures อื่นใน paper เดียวกัน การทดลองเปรียบเทียบ No-Profile, One-Profile และ All-Profile เพื่อดูว่าจำนวน profile ส่งผลต่อ caption generation อย่างไร

\subsection*{ผลลัพธ์สำคัญ}
\begin{itemize}
\item การใช้ profile information ทำให้ caption ใกล้ reference caption มากขึ้นอย่างสม่ำเสมอ
\item All-Profile มักให้คะแนนดีที่สุด แต่ gain จาก No-Profile ไป One-Profile ใหญ่กว่าการเพิ่ม profile หลังจากนั้น แปลว่าการเลือก context ที่ relevant สำคัญกว่าการใส่ context เยอะเสมอ
\item ablation พบว่า profile caption มีผลมากที่สุด แต่ profile image ก็มีประโยชน์มากกว่า figure-mentioning paragraph ในบาง setting
\item BERTScore และ ROUGE-L distribution ระหว่าง target กับ profile captions ชี้ว่า profile มี semantic relatedness สูง แต่ lexical overlap ต่ำ เหมาะกับการเป็น style/context signal
\end{itemize}

\subsection*{ข้อจำกัด}
\begin{itemize}
\item profile ถูกจำกัดสูงสุด 3 รูป ไม่ใช่ full-document context ทั้ง paper
\item task ยังใช้ context ที่ถูก curate มาแล้วจาก dataset ไม่ได้เริ่มจาก raw PDF parsing
\item evaluation ยังพึ่ง BLEU/ROUGE/BERTScore ซึ่งอาจไม่วัด faithfulness ต่อภาพหรือความถูกต้องเชิงวิชาการครบถ้วน
\item การใกล้ author-written caption ไม่ได้แปลว่า caption ดีกว่าเสมอ เพราะ original caption อาจไม่ informative
\end{itemize}

\subsection*{ความเกี่ยวข้องกับ Thesis}
Paper นี้เป็นฐานที่ดีสำหรับ research gap ของ thesis เพราะพิสูจน์แล้วว่า context จาก paper เดียวกันช่วย scientific figure captioning ได้ แต่ยังเปิดช่องว่างสำคัญคือ context ยังเป็น profile figures ที่ curated แล้ว ไม่ใช่ context retrieval จาก academic PDF ทั้งฉบับ Thesis สามารถต่อยอดเป็น \textbf{multi-level document context retrieval} เช่น mention paragraph, section context, abstract, nearby text, OCR และ cross-reference แล้วนำไปใช้กับ VLM/RAG caption generation

\newpage

\section{Leveraging Author-Specific Context for Scientific Figure Caption Generation: 3rd SciCap Challenge}

\noindent
\begin{tabular}{@{}ll@{}}
\textbf{arXiv} & \href{https://arxiv.org/abs/2510.07993}{2510.07993} \\
\textbf{ผู้เขียน} & Watcharapong Timklaypachara, Monrada Chiewhawan, Nopporn Lekuthai, Titipat Achakulvisut \\
\textbf{Keyword} & Context-Aware Caption Generation, Author-Specific Context, Prompt Optimization \\
\textbf{ไฟล์ใน Repo} & papers/9-AuthorContextSciCap\_2025.* \\
\end{tabular}

\subsection*{ปัญหาที่ Paper แก้}
งานนี้เริ่มจาก observation ว่า zero-shot prompting ที่ใช้ paragraph information โดยตรงมักสร้าง caption ที่ noisy มี irrelevant content และ hallucination โดยเฉพาะเมื่อ paragraph ยาวหรือมีข้อมูลที่ไม่ได้เกี่ยวกับ figure โดยตรง Scientific figure caption จึงต้องแก้สองโจทย์พร้อมกัน: ต้อง ground กับ content จริง และต้องมี style/structure ที่ใกล้กับ paper ต้นฉบับ

\subsection*{Contribution สำคัญ}
\begin{enumerate}
\item เสนอ two-stage caption generation pipeline สำหรับ 3rd SciCap Challenge
\item Stage 1 ทำ content-grounded caption generation โดยใช้ context filtering และ category-focused prompt optimization
\item Stage 2 ทำ profile-informed stylistic refinement โดยใช้ few-shot prompting จาก profile figures
\item ใช้ DSPy MIPROv2 และ SIMBA เพื่อ optimize prompt ตาม paper category
\item ใช้ Gemini-2.5-flash เป็น reranker/candidate selector เพื่อเลือก caption ที่เหมาะกว่า
\end{enumerate}

\subsection*{Method}
Pipeline แบ่งเป็นสอง stage:

\paragraph{Stage 1: Content-Grounded Caption Generation}
ระบบกรอง paragraph เพื่อเอาประโยคที่เกี่ยวกับ figure มากขึ้น จากนั้นใช้ category-focused prompt templates เพราะ paper ต่าง field มี terminology และ pattern ของ caption ต่างกัน MIPROv2 สร้าง instruction-example pairs และ optimize ด้วย ROUGE-L precision ส่วน SIMBA เพิ่ม feedback-driven optimization เพื่อหา reasoning path และปรับ rule สำหรับ problematic cases

\paragraph{Stage 2: Profile-Informed Stylistic Refinement}
หลังได้ caption จาก Stage 1 ระบบใช้ profile figures จาก LaMP-Cap เป็น few-shot examples เพื่อปรับ style ให้ใกล้กับ paper ต้นฉบับ และบังคับความยาวประมาณ \(\pm 15\%\) เพื่อคุม conciseness จุดนี้สำคัญเพราะ Stage 1 อาจได้ caption ที่มี content coverage ดี แต่ยัง verbose หรือไม่เหมือน writing style ของผู้เขียน

\subsection*{ผลลัพธ์สำคัญ}
\begin{highlight}
\textbf{Key numbers}
\begin{itemize}[nosep]
\item Filtering + prompt optimization เพิ่ม ROUGE-1 recall จาก 0.4867 เป็น 0.5443
\item ROUGE-2 recall เพิ่มจาก 0.2223 เป็น 0.2572
\item Category-focused prompting เพิ่ม ROUGE-1 recall ประมาณ +8.3\% และลด BLEU-4 น้อยกว่า non-category-focused prompting
\item Gemini reranking เพิ่ม ROUGE precision ประมาณ 3.46--4.43\%
\item Profile-informed refinement เพิ่ม BLEU ประมาณ 40--48\% และ precision-based ROUGE ประมาณ 25--27\%
\end{itemize}
\end{highlight}

\subsection*{ข้อจำกัด}
\begin{itemize}
\item pipeline มีหลาย component ทำให้ cost, reproducibility และ implementation complexity สูง
\item category prompt optimization ยังไม่ใช่ retrieval จาก document structure โดยตรง
\item การใช้ BLEU/ROUGE อาจให้คะแนนกับ lexical/style overlap มากกว่าความถูกต้องของข้อมูลในภาพ
\item ใช้ profile context เป็นหลัก ยังไม่ได้ใช้ layout-aware PDF context เช่น section, table, caption proximity หรือ cross-reference อย่างเป็นระบบ
\end{itemize}

\subsection*{ความเกี่ยวข้องกับ Thesis}
Paper นี้ช่วยยืนยันว่า \textbf{context filtering ก่อน generation} เป็นขั้นตอนจำเป็น ไม่ควรส่ง paragraph ทั้งหมดให้โมเดลโดยไม่กรอง Thesis สามารถนำแนวคิดนี้ไปออกแบบ retrieval module ที่เลือก context จากเอกสารทั้งฉบับแบบอัตโนมัติ แทนที่จะพึ่ง category prompt และ profile-only refinement Research gap ที่ชัดคือการสร้าง pipeline ที่เริ่มจาก PDF จริง แล้วทำ layout-aware context extraction ก่อนส่งเข้า VLM

\newpage

\section{Personalized Scientific Figure Caption Generation: An Empirical Study on Author-Specific Writing Style Transfer}

\noindent
\begin{tabular}{@{}ll@{}}
\textbf{arXiv} & \href{https://arxiv.org/abs/2509.25817}{2509.25817} \\
\textbf{ผู้เขียน} & Jaeyoung Kim, Jongho Lee, Hongjun Choi, Sion Jang \\
\textbf{Keyword} & Scientific Figure Captioning, Quality-Aware Generation, Qwen2.5-VL \\
\textbf{ไฟล์ใน Repo} & papers/10-PersonalizedSciFigCap\_2025.* \\
\end{tabular}

\subsection*{ปัญหาที่ Paper แก้}
Paper นี้ศึกษาว่า profile data ช่วย personalized scientific figure captioning ได้จริงแค่ไหน และชี้ปัญหาสำคัญว่า metric แบบ BLEU/ROUGE อาจทำให้โมเดลเรียนรู้การเลียนแบบ author style หรือ lexical pattern มากกว่าการสร้าง caption ที่ informative และ factual งานก่อนหน้าระบุว่า caption ใน arXiv จำนวนมากมีคุณภาพต่ำ ดังนั้นการ fine-tune จาก caption เหล่านี้โดยตรงอาจทำให้โมเดลเรียนรู้ caption ที่ไม่ดีตามไปด้วย

\subsection*{Contribution สำคัญ}
\begin{enumerate}
\item สร้าง quality evaluator โดยใช้ GPT-4.1 สร้าง synthetic labels ระดับ 1--6 แล้ว fine-tune Qwen-2.5-VL-3B
\item fine-tune Qwen-2.5-VL-7B เป็น personalized caption generator จาก LaMP-Cap
\item วิเคราะห์ input components ได้แก่ figure, paragraph, mention, OCR และ profile data
\item เสนอ quality-aware caption generation เพื่อควบคุม trade-off ระหว่าง personalization และ caption quality
\item ชี้ limitation ของ BLEU/ROUGE เมื่อ reference caption เองอาจมีคุณภาพต่ำ
\end{enumerate}

\subsection*{Method}
ระบบมีสองส่วนหลัก ส่วนแรกคือ \(f_{quality}\) สำหรับประเมิน caption quality โดยใช้ figure-caption pairs และให้คะแนน 1 ถึง 6 ตามคุณภาพ ส่วนนี้ fine-tune จาก Qwen-2.5-VL-3B ด้วย synthetic labels จาก GPT-4.1 จำนวน 3,000 samples และ validate บน 200 samples

ส่วนที่สองคือ \(g_{caption}\) ใช้ Qwen-2.5-VL-7B เป็น generator โดย input มี target figure, explanatory paragraphs, textual mentions, OCR text และ profile data จาก paper เดียวกัน สำหรับ supervised fine-tuning ใช้เฉพาะ samples ที่ quality score \(\geq 3\) เพื่อเลี่ยง low-quality captions

\begin{highlight}
\textbf{ตัวเลขจาก quality evaluator}
\begin{itemize}[nosep]
\item Spearman rank correlation = 0.759
\item Quadratic Weighted Kappa = 0.754
\item training samples สำหรับ quality assessment = 3,000
\item validation samples = 200
\end{itemize}
\end{highlight}

\subsection*{ผลลัพธ์สำคัญ}
\begin{itemize}
\item Fine-tuned Qwen2.5-VL-7B ทำคะแนนบน full test ดีกว่า Gemini-2.5-flash: BLEU-4 = 0.304 เทียบกับ 0.278 และ ROUGE-L = 0.555 เทียบกับ 0.540
\item จำนวน profile ที่มากขึ้นช่วยเพิ่ม BLEU/ROUGE อย่างสม่ำเสมอ โดยตัวอย่างที่มี 3 profile figures ได้คะแนนสูงสุดในหลาย metric
\item Ablation พบว่าเพิ่ม paragraphs จาก baseline F+C ทำให้ BLEU-4 เพิ่มจาก 0.123 เป็น 0.295 และ ROUGE-L เพิ่มจาก 0.355 เป็น 0.545
\item mentions และ OCR เพิ่มผลแบบ incremental จน full configuration ได้ BLEU-4 = 0.304 และ ROUGE-L = 0.555
\item quality-aware generation แบบ Forced-Q6 ทำให้ quality score สูงขึ้น แต่ ROUGE-L อาจลดในกรณี reference caption คุณภาพต่ำ
\end{itemize}

\subsection*{ข้อจำกัด}
\begin{itemize}
\item quality labels มาจาก GPT-4.1 จึงอาจสะท้อน bias ของ judge model
\item optimization เพื่อ style matching อาจลด informativeness ของ caption
\item evaluation หลักยังพึ่ง lexical overlap metrics ทำให้ caption ที่ดีกว่าแต่ wording ต่างจาก reference อาจได้คะแนนต่ำ
\item ยังไม่ได้เป็น pipeline จาก raw academic PDF เพราะใช้ LaMP-Cap fields ที่เตรียมไว้แล้ว
\end{itemize}

\subsection*{ความเกี่ยวข้องกับ Thesis}
Paper นี้ให้ insight สำคัญมากสำหรับ thesis: context components อย่าง paragraph, mention และ OCR มีผลชัดเจนต่อ caption quality แต่การประเมินด้วย BLEU/ROUGE เพียงอย่างเดียวอาจไม่พอ Thesis ควรแยก evaluation อย่างน้อย 4 มิติ คือ visual faithfulness, context relevance, academic informativeness และ lexical/style similarity นอกจากนี้แนวทาง quality-aware generation สามารถต่อยอดเป็น post-generation verification หรือ refinement module เพื่อลด hallucination

\newpage

\section*{Research Gap ที่สรุปจากทั้ง 3 Paper}

\begin{enumerate}
\item \textbf{Context Integration Gap:} งานปี 2025 ใช้ context มากขึ้น แต่ยังมักเป็น profile figures หรือ fields ที่ curated แล้ว ยังไม่ใช่ retrieval จาก PDF ทั้งฉบับ
\item \textbf{End-to-End Pipeline Gap:} ยังขาดระบบครบวงจรจาก PDF parsing, layout analysis, figure extraction, context retrieval ไปจนถึง caption generation
\item \textbf{Evaluation Gap:} BLEU/ROUGE/BERTScore วัด similarity กับ reference แต่ไม่พอสำหรับ faithfulness, factuality และ informativeness
\item \textbf{Quality-Style Trade-off:} การทำให้ caption เหมือน author-written caption อาจไม่เท่ากับการทำให้ caption ดีขึ้น เพราะ reference caption บางส่วนมีคุณภาพต่ำ
\item \textbf{Retrieval Quality Gap:} ยังไม่มีคำตอบชัดว่าควรเลือก context แบบใด ระหว่าง profile figures, mention paragraph, section context, OCR, nearby text หรือ full-document RAG
\end{enumerate}

\begin{highlight}
\textbf{ข้อเสนอสำหรับ thesis direction:}
พัฒนา Hybrid RAG + VLM pipeline ที่เริ่มจาก academic PDF จริง ทำ layout-aware document parsing แล้วดึง multi-level context สำหรับ target figure ก่อนสร้าง caption พร้อม evaluation ที่แยก visual faithfulness, context relevance และ academic informativeness ออกจาก lexical overlap
\end{highlight}

\vspace{16pt}
\noindent\rule{\textwidth}{1pt}
\noindent{\color{muted}\small สร้างด้วย LaTeX skill · อ้างอิงจาก papers/8--10 และ docs/synthesis analysis · \today}

\end{document}
